\documentclass[12pt]{article}
\usepackage[utf8]{inputenc}
\usepackage[T1]{fontenc}
\usepackage{hyperref}
\title{Wakon\'a Language Resources for K-12 Curriculum Development}
\author{Compiled by Claude Sonnet 5 (claude-sonnet-5) for Andr\'as Kornai}
\date{}
\begin{document}
\maketitle

\section{Introduction}
Wakon\'a (also Acon\~a; ISO 639-3 \texttt{waf}) was spoken in Alagoas, Brazil. It is unclassified and no longer used as a first language. Wikipedia describes it as extinct, while Ethnologue lists it as dormant. The only known data is a very short word list of about nine items, collected in 1961 and published by Meader in 1978. The Acon\~a are one of the ethnic components of the Xukuru-Kariri people.

\section{Primary Educational Resources}
None were found. A K-12 curriculum in the language is not feasible from the attested data.

\section{Linguistic and Reference Materials}
\begin{itemize}
\item Wikipedia article on the Wakon\'a language: \url{https://en.wikipedia.org/wiki/Wakon%C3%A1_language}
\item Ethnologue entry: \url{https://www.ethnologue.com/language/waf}
\item A UFRJ repository record seen in the results (content not checked): \url{https://bdor.sibi.ufrj.br/handle/doc/103}
\end{itemize}

\section{Religious and Community Resources}
None were found specific to Wakon\'a. The Xukuru-Kariri people are the related community.

\section{Digital and Audio Resources}
None were found.

\section{Curriculum Development Context}
The language is effectively unattested. Any educational use would be limited to cultural or historical awareness, in cooperation with the Xukuru-Kariri community, with Portuguese as the language of instruction.
\end{document}